\chapter{Polar alignment and perimeter growth}\label{sec:perimeter-polar-growth}
The reduced-boundary definition contains only a limiting average of the distributional derivative. We develop its quantitative consequences before using any rectifiability theorem. Throughout this section $n\geq2$, $E$ has locally finite perimeter in an open set $U\subset\mathbb R^n$, and
\[
 \mu=|D\chi_E|,\qquad D\chi_E=-\nu\mu,\qquad |\nu|=1\quad\mu\text{-almost everywhere}.
\]
Here $\nu$ is a measurable outward polar field, and $\nu_0=\nu_E(x)$ is its limiting average at $x\in\partial^*E$. This is the outward convention of Definition \ref{def:reduced-boundary}. The source argument is \cite{Sim18}, Chapter 3, Theorem 4.4, equations (1)--(6), printed pp.~97--98. None of the results below uses the De Giorgi structure theorem.

\begin{lemma}[The squared polar excess is an average defect]\label{lem:perimeter-polar-excess}
For $x\in\partial^*E$ and all sufficiently small $r>0$, define
\[
 a(r)=\frac{1}{\mu(B_r(x))}\int_{B_r(x)}\nu\cdot\nu_0\,d\mu,
 \qquad \varepsilon(r)=\frac{1}{\mu(B_r(x))}
                      \int_{B_r(x)}|\nu-\nu_0|^2\,d\mu.
\]
Then $\varepsilon(r)=2(1-a(r))\to0$. In particular,
\[
 \int_{B_r(x)}|\nu-\nu_0|\,d\mu
 \leq\sqrt{\varepsilon(r)}\,\mu(B_r(x)),
 \qquad a(r)\geq\tfrac12
\]
for every sufficiently small radius. This expands the polar-component estimates in \cite{Sim18}, equations (1)--(2).
\end{lemma}
\begin{proof}
The denominators are positive because $x$ belongs to the support of $\mu$. Both polar vectors have unit length, so $|\nu-\nu_0|^2=2-2\nu\cdot\nu_0$ almost everywhere. Integrating gives the identity. The reduced-boundary definition says that the average of $\nu$ tends to $\nu_0$, so $a(r)\to1$. Cauchy--Schwarz with respect to the finite measure on the ball gives the first estimate. The last assertion follows from the limit, for all sufficiently small $r$, without selecting a sequence.
\end{proof}

\begin{corollary}[Tangential and negative polar components disappear]\label{cor:perimeter-polar-components}
Let $P=I-\nu_0\otimes\nu_0$ be projection onto $\nu_0^\perp$, and let $s_- =\max\{-s,0\}$. Under Lemma \ref{lem:perimeter-polar-excess},
\[
 \int_{B_r(x)}|P\nu|\,d\mu\leq\sqrt{\varepsilon(r)}\mu(B_r(x)),
 \qquad
 \int_{B_r(x)}(\nu\cdot\nu_0)_-\,d\mu
 \leq\tfrac12\varepsilon(r)\mu(B_r(x)).
\]
The normal component need not be nonnegative at every point of a small ball. These estimates, rather than pointwise sign assertions, pass to blow-up limits; compare \cite{Sim18}, pp.~97--98.
\end{corollary}
\begin{proof}
Since $P\nu_0=0$, one has $|P\nu|\leq|\nu-\nu_0|$. Apply the preceding lemma. For $-1\leq s\leq1$, $s_-\leq1-s$. With $s=\nu\cdot\nu_0$, integration gives the second estimate from the exact excess identity.
\end{proof}

\begin{lemma}[Radial phase volumes are absolutely continuous]\label{lem:radial-phase-volume}
If $B_R(x)\Subset U$, set $V(r)=|E\cap B_r(x)|$ and $W(r)=|B_r(x)\setminus E|$ for $0\leq r<R$. These functions are absolutely continuous, and for almost every $r$,
\[
 V'(r)=\int_{\partial B_r(x)}\chi_E\,d\mathcal H^{n-1},\qquad
 W'(r)=n\omega_n r^{n-1}-V'(r),
 \qquad 0\leq V'(r),W'(r)\leq n\omega_n r^{n-1}.
\]
Representatives of $E$ differing on a volume-null set give the same identities for almost every radius. This is the radial differentiation used in \cite{Sim18}, equations (4)--(6).
\end{lemma}
\begin{proof}
Translate $x$ to zero and use polar coordinates and Tonelli:
\[
 V(r)=\int_0^r s^{n-1}\int_{S^{n-1}}\chi_E(s\theta)\,d\mathcal H^{n-1}(\theta)\,ds.
\]
The integrand is measurable and bounded by $n\omega_n s^{n-1}$. The fundamental theorem for Lebesgue integrals proves absolute continuity and the derivative formula. The sphere integral equals the displayed derivative after dilation. Subtract $V$ from $\omega_n r^n$ to obtain the claims for $W$. Fubini makes a volume-null modification invisible on almost every sphere.
\end{proof}

\begin{lemma}[A radial flux bound without a smooth boundary]\label{lem:perimeter-radial-flux}
For any unit vector $e$ and almost every $r\in(0,R)$ with $B_R(x)\Subset U$,
\[
 \left|\int_{B_r(x)}\nu\cdot e\,d\mu\right|\leq V'(r).
\]
The corresponding bound for $U\setminus E$ has $-\nu$ in place of $\nu$ and $W'(r)$ in place of $V'(r)$. This is a distributional version of the flux estimate in \cite{Sim18}, equations (3)--(5); no boundary trace on the sphere is required.
\end{lemma}
\begin{proof}
Choose a smooth nonincreasing function $h$ on $\mathbb R$, equal to one on $(-\infty,0]$ and zero on $[1,\infty)$, with $\int_0^1-h'(s)\,ds=1$. Put $\phi_\delta(y)=h((|y-x|-r)/\delta)$ for $0<\delta<R-r$. It is smooth, including near $x$, because it is constant there. The identity $D\chi_E=-\nu\mu$ gives
\[
 \int\phi_\delta\nu\cdot e\,d\mu
 =\int_E e\cdot D\phi_\delta\,dy,
 \qquad
 \left|\int_E e\cdot D\phi_\delta\right|
 \leq\int_0^1(-h'(s))V'(r+\delta s)\,ds.
\]
For almost every $r$, the right side tends to $V'(r)$ by Lebesgue differentiation applied to the integrable radial derivative from Lemma \ref{lem:radial-phase-volume}. The exceptional radii where $\mu(\partial B_r(x))>0$ are at most countable, since these spheres are disjoint and $\mu$ is finite on each smaller compact ball. At every remaining radius, dominated convergence gives convergence of the left integral to its open-ball integral. Taking absolute values proves the result. The complement satisfies $D\chi_{U\setminus E}=-D\chi_E$ within $U$, so the same proof applies to it.
\end{proof}

\begin{proposition}[Perimeter growth at a reduced-boundary point]\label{prop:reduced-boundary-perimeter-growth}
For $x\in\partial^*E$ there is $r_0>0$ such that for almost every $0<r<r_0$,
\[
 \mu(B_r(x))\leq2V'(r),\qquad
 \mu(B_r(x))\leq2W'(r).
\]
Moreover, at every $0<r<r_0$,
\[
 \mu(B_r(x))\leq2n\omega_n r^{n-1}.
\]
The radius $r_0$ can depend on $E$ and $x$, but the upper-growth constant depends only on $n$. This is the compactness bound in \cite{Sim18}, equation (6), made valid at all small radii.
\end{proposition}
\begin{proof}
Choose $r_0$ so that Lemma \ref{lem:perimeter-polar-excess} gives $a(r)\geq1/2$ for all $r<r_0$, and the closed ball of radius $r_0$ lies in $U$. The flux lemma with $e=\nu_0$ gives $\mu(B_r)/2\leq\int_{B_r}\nu\cdot\nu_0\,d\mu\leq V'(r)$ at good radii. For the complement take outward limiting vector $-\nu_0$; the polar dot product is unchanged, and the same reasoning gives the bound with $W'$. Lemma \ref{lem:radial-phase-volume} bounds either derivative by $n\omega_n r^{n-1}$. Finally, for any radius $r<r_0$ choose good radii $s_j\downarrow r$ with $s_j<r_0$. Monotonicity gives $\mu(B_r)\leq\mu(B_{s_j})\leq2n\omega_n s_j^{n-1}$; taking the limit proves the all-radius bound without assuming the sphere has zero perimeter mass.
\end{proof}

\begin{lemma}[Scaling the distributional derivative and its variation]\label{lem:perimeter-blowup-scaling}
For $x\in U$ and $r>0$, write $F_r(y)=x+ry$, $U_{x,r}=F_r^{-1}(U)$ and $E_{x,r}=F_r^{-1}(E)$. Then in $U_{x,r}$,
\[
 D\chi_{E_{x,r}}=r^{1-n}(F_r^{-1})_\#D\chi_E,
 \qquad |D\chi_{E_{x,r}}|=r^{1-n}(F_r^{-1})_\#\mu.
\]
Thus the outward polar field is $\nu(x+ry)$ almost everywhere for the rescaled perimeter measure. This supplies the measure scaling used in \cite{Sim18}, equation (8).
\end{lemma}
\begin{proof}
For a compactly supported smooth vector field $X$ on $U_{x,r}$, change variables in $\int\chi_{E_{x,r}}\operatorname{div}X$. The chain rule contributes $r$ when the test is expressed as $X((z-x)/r)$ in the original variable, whereas volume contributes $r^{-n}$. The distributional integration-by-parts identity gives the first formula. A bijective measurable change of variables preserves the norm of the polar vector; its pushforward is still unit almost everywhere. The vector-measure polar representation therefore gives the second formula. This is scalar dilation in fixed Euclidean coordinates; no tangent-plane measure formula is being assumed.
\end{proof}
